\chapter{Methodology}
\label{ch:methodology}

\section{The kernel ladder}

The point of the ladder is attribution. Three force kernels differ by exactly
one technique each, so the difference between two rungs can be assigned to that
technique and nothing else. All three compute the same quantity from the same
structure of arrays layout, and a test asserts they agree with a double
precision host reference to better than one part in ten thousand, so a
performance difference cannot be a difference in what is being computed.

\paragraph{Rung one, naive.} Every thread walks the whole particle array in
device memory. This is not a straw man built to lose. All threads in a warp read
the same index on the same cycle, so the load is a broadcast and the line stays
resident for the whole warp. What it does not do is share that line between the
warps of a block, so the same data crosses the cache boundary once per resident
block rather than once per tile.

\paragraph{Rung two, tiled.} The block cooperatively stages a tile of source
particles in shared memory and every thread of the block consumes it, following
Nyland, Harris and Prins \cite{nyland2007}. The tile width is thirty two by
default. The cost is two barriers per tile, so narrow tiles pay the barrier
more often per unit of arithmetic; Section~\ref{sec:tilewidth} measures that
trade rather than arguing about it.

\paragraph{Rung three, fused.} The tiled force evaluation and the integrator
update happen in one launch. This removes the read back of the acceleration
array, $12N$ bytes per step, plus one kernel launch per step. The acceleration
is still written, because the energy diagnostic needs it.

Fusion has a correctness constraint that is worth stating because it is easy to
get wrong. There is no grid wide barrier inside a launch. If the kernel updated
positions in place, a block that had finished its force sum could overwrite a
position that another block was still reading, and the result would depend on
scheduling. Positions are therefore double buffered: read from one array,
written to the other, swapped afterwards. This costs $12N$ bytes of memory and
no bandwidth at all, because the write has to happen either way.

\section{Timing protocol}

Three things ruin a GPU timing measurement, and the protocol addresses each.

The first launch of a kernel pays for module loading and for the card ramping
up its clocks, so every measurement is preceded by warm up steps that are not
timed.

The host can queue work far faster than the device retires it, so timing on the
host measures the host. All timings use CUDA events recorded on the same stream
as the work.

Run length varies enormously across the sweep, from a fraction of a millisecond
to nearly a hundred. Each point is therefore calibrated first with a single
timed step, and the run length is chosen to make each trial last about
\SI{120}{\milli\second}.

Each point runs three such trials and the fastest is kept. The minimum rather
than the mean, because everything that perturbs a measurement, whether thermal
throttling or another process on the display, makes it slower and never faster.
The minimum is the estimate least contaminated by whatever else the machine was
doing.

One caveat is worth recording. Absolute timings depend on the thermal state of
the card. A tile width sweep run immediately after the full benchmark sweep
produced timings about fifteen percent higher than the same sweep run from
cold, while the relative ordering and the percentage differences between tile
widths were reproduced to within a tenth of a percentage point. The comparisons
in this report are therefore all made within a single session, and the
conclusions rest on relative differences rather than on absolute
milliseconds.

\section{Validation design}

Each gate is chosen to fail for a specific class of defect rather than to
produce a reassuring number.

\paragraph{Two body energy drift.} A circular orbit integrated for ten thousand
steps must conserve total energy to better than one part in a thousand. This is
the sharpest single statement that can be made about an integrator, and it also
catches an inconsistent softened potential and a mis-derived initial orbital
speed. The initial speed is computed from the softened force law, so the orbit
is genuinely circular rather than a slowly precessing ellipse.

\paragraph{Momentum conservation.} The internal forces of an isolated system
cancel exactly in the mathematics. On the device they cancel only to rounding,
because each thread accumulates its own sum in single precision in its own
order. The gate bounds the residual against what single precision can explain,
$\epsilon_{\mathrm{f32}} \sqrt{N} n_{\mathrm{steps}}$, rather than against an
arbitrary constant.

\paragraph{Force antisymmetry.} Newton's third law checked pair by pair against
an independent double precision host implementation, written as literally as
the force law can be written, with no tiling and no reciprocal square root. Its
only job is to be obviously correct.

\paragraph{Plummer radial density.} A chi squared test of the sampled radial
distribution against the analytic cumulative mass profile. Shells are chosen to
hold equal expected counts, so no bin is too sparse for the statistic to be
valid. A sampler with a subtly wrong exponent produces a blob that looks
perfectly reasonable and fails this immediately. The test is itself checked by
running it against a deliberately wrong scale radius, which it must reject.

\paragraph{SPH kernel normalisation.} Each smoothing kernel is integrated over
its support by composite Simpson quadrature and must come to one within
$10^{-3}$. A kernel whose normalisation is off by a few percent produces a
density field wrong by the same few percent everywhere, and the fluid still
looks like a fluid.

\paragraph{Neighbour search against brute force.} The density field from the
hashed neighbour search is compared against an $O(N^2)$ host sum over the same
positions. This is what would catch a dropped or double counted neighbour, and
it is the direct test of the cell key defence described in
Chapter~\ref{ch:background}.

\section{Reproducibility}

Every configuration carries an explicit seed and the initial conditions are
generated on the host from a written out uniform mapping rather than from a
standard library distribution, whose sequence is not specified across
implementations. Each benchmark row records the seed, the commit, the CUDA
runtime and driver versions, the fast math setting and a timestamp. Snapshots
are written to a temporary file and renamed into place, so an interrupted run
never leaves a half written snapshot for a later resume to pick up, and a test
asserts that stopping and resuming reproduces a straight through run bit for
bit.
